% Team report source. Edit this section in a feature branch and request peer review.
\needspace{4\baselineskip}\section{Task 4: Optimized design for Workload B}
As suggested in the project description, the sequence of optimizations was progressively conducted. In general, each optimization was selected based on the measurements of the previous experiment, we selected the optimizations based on the workload weights and the performance bottlenecks observed in previous experiments. Within the allowed tuning scope we changed the data distribution, the schema, the indexes and the transaction code.\par
\needspace{4\baselineskip}\subsection{\texttt{expt\_\allowbreak{}b1}}
T7 makes up 30\% of workload B, and based on the vanilla baseline schema, each of its four \texttt{SELECT} queries will scan all 32 shards of \texttt{inventory} because they filter on \texttt{product\_\allowbreak{}id}. In addition, T4 requires a repartition join between \texttt{products} and \texttt{inventory} on \texttt{product\_\allowbreak{}id}. Distributing \texttt{inventory} on \texttt{product\_\allowbreak{}id}, in the same co-location group as \texttt{products}, optimises both. The expected outcome is every T7 \texttt{SELECT} routes to one shard, and the T4 join becomes local to one shard.\par
\needspace{4\baselineskip}\subsubsection*{Key Changes}
Change the definition of \texttt{inventory}  so the primary key becomes \texttt{(product\_\allowbreak{}id, inventory\_\allowbreak{}id)};\par
\texttt{inventory} to be distributed on \texttt{product\_\allowbreak{}id} and \texttt{colocate\_\allowbreak{}with => 'products'}.\par
T7's two updates retain the selected inventory IDs and add their selected product IDs as routing predicates. The supplied fixture happens to use inventory\_id equal to product\_id, but the SQL looks up actual inventory IDs and does not rely on that equality.\par
\needspace{4\baselineskip}\subsubsection*{Benchmark Results and Analysis}
\needspace{10\baselineskip}
{\fontsize{10.5}{12.5}\selectfont\setlength{\tabcolsep}{2mm}\renewcommand{\arraystretch}{1.15}
\begin{longtable}{>{\raggedright\arraybackslash}p{30.0mm}>{\raggedright\arraybackslash}p{30.0mm}>{\raggedright\arraybackslash}p{30.0mm}>{\raggedright\arraybackslash}p{30.0mm}>{\raggedright\arraybackslash}p{30.0mm}}\toprule
 & \textbf{TPS} & \textbf{p50 Latency (ms)} & \textbf{p95 Latency (ms)} & \textbf{p99 Latency (ms)} \\
\midrule\endhead
T7 & 0.44 \ensuremath{\rightarrow} 0.73 & 1,517 \ensuremath{\rightarrow} 275 & 17,736 \ensuremath{\rightarrow} 2,164 & 26,139 \ensuremath{\rightarrow} 3,465 \\
T4 & 0.21 \ensuremath{\rightarrow} 0.12 & 74,221 \ensuremath{\rightarrow} 53 & 84,350 \ensuremath{\rightarrow} 216 & 87,024 \ensuremath{\rightarrow} 381 \\
Overall & 3.19 \ensuremath{\rightarrow} 2.26 & 35 \ensuremath{\rightarrow} 249 & 133,192 \ensuremath{\rightarrow} 18,335 & 160,032 \ensuremath{\rightarrow} 926,499 \\
\bottomrule\end{longtable}}
{\itshape Table 5.1.1: Comparing expt\_b1 against expt\_b0 at 64 clients, on T4, T7  --  flagged raw observations; aborted phases exclude a clean comparison.}\par
The flagged c64 observations show lower T7/T4 latencies and fewer recorded serialization failures, consistent with the routing hypothesis. Both b0 and b1 c64 warmups aborted, and their allocations differ (xcnf versus xcnc); these ratios do not establish a clean matched causal effect.\par
The raw c64 throughput falls from 3.19 to 2.26 TPS, but aborted warmups and different hardware prevent a clean comparison. T3 remains a candidate bottleneck in the recorded time shares.\par
A T7-only supplementary comparison uses \texttt{expt\_b201}/job 989344 and baseline \texttt{expt\_b0}/job 988876. At one client, baseline b0 records 5.188 TPS and 192.786 ms mean latency, while b201 records 334.518 TPS and 2.984 ms. The baseline uses xcnc34-38 while the optimized run uses xcnf10-14, so the full difference cannot be attributed to the routing change alone. T7 is 100\% of this supplementary workload; these observations do not replace the original B mix.\par
\needspace{4\baselineskip}\subsection{\texttt{expt\_\allowbreak{}b2}}
T3 has accounted for most of the time during the benchmarking runs up till this point. Its join of \texttt{orders} and \texttt{order\_\allowbreak{}items} on \texttt{order\_\allowbreak{}id} was a repartition join and required shifting rows to different workers as \texttt{order\_\allowbreak{}items} was distributed on \texttt{order\_\allowbreak{}item\_\allowbreak{}id} rather than \texttt{order\_\allowbreak{}id}.\par
\needspace{4\baselineskip}\subsubsection*{Key Changes}
Change the primary key of \texttt{order\_\allowbreak{}items} to be \texttt{(order\_\allowbreak{}id, order\_\allowbreak{}item\_\allowbreak{}id)}.\par
\texttt{order\_\allowbreak{}items} is distributed on \texttt{order\_\allowbreak{}id}  and is co-located with \texttt{orders} via \texttt{colocate\_\allowbreak{}with => 'orders'}.\par
\needspace{4\baselineskip}\subsubsection*{Benchmark Results}
\needspace{10\baselineskip}
{\fontsize{10.5}{12.5}\selectfont\setlength{\tabcolsep}{2mm}\renewcommand{\arraystretch}{1.15}
\begin{longtable}{>{\raggedright\arraybackslash}p{30.0mm}>{\raggedright\arraybackslash}p{30.0mm}>{\raggedright\arraybackslash}p{30.0mm}>{\raggedright\arraybackslash}p{30.0mm}>{\raggedright\arraybackslash}p{30.0mm}}\toprule
 & \textbf{TPS} & \textbf{p50 Latency (ms)} & \textbf{p95 Latency (ms)} & \textbf{p99 Latency (ms)} \\
\midrule\endhead
T3 & 0.11 \ensuremath{\rightarrow} 2.37 & 885,875 \ensuremath{\rightarrow} 1,381 & 941,152 \ensuremath{\rightarrow} 2,139 & 956,702 \ensuremath{\rightarrow} 2,422 \\
T6 & 0.20 \ensuremath{\rightarrow} 4.61 & 15,967 \ensuremath{\rightarrow} 12,443 & 17,829 \ensuremath{\rightarrow} 13,782 & 18,483 \ensuremath{\rightarrow} 14,275 \\
Overall & 2.26 \ensuremath{\rightarrow} 46.04 & 249 \ensuremath{\rightarrow} 68 & 18,335 \ensuremath{\rightarrow} 12,444 & 926,499 \ensuremath{\rightarrow} 13,510 \\
\bottomrule\end{longtable}}
{\itshape Table 5.2.1: Comparing expt\_b2 against expt\_b1 on T3, T6 at 64 clients  --  flagged raw observations; aborted phases exclude a clean comparison.}\par
The c64 raw observations suggest much lower T3 latency after order/item colocation, but b1's warmup and b2's measurement aborted. The apparent 20-fold TPS ratio is diagnostic, not an eligible matched speedup. b2 has complete zero-failure phases at 1/4/8/16 clients.\par
\needspace{10\baselineskip}
{\fontsize{10.5}{12.5}\selectfont\setlength{\tabcolsep}{2mm}\renewcommand{\arraystretch}{1.15}
\begin{longtable}{>{\raggedright\arraybackslash}p{30.0mm}>{\raggedright\arraybackslash}p{30.0mm}>{\raggedright\arraybackslash}p{30.0mm}>{\raggedright\arraybackslash}p{30.0mm}>{\raggedright\arraybackslash}p{30.0mm}}\toprule
\textbf{Txn} & \textbf{Avg latency, 1 client (ms)} & \textbf{Share of time, 1 client} & \textbf{Avg latency, 64 clients} & \textbf{Share of time, 64 clients} \\
\midrule\endhead
T3 & 10.9 & 1.7\% & 1,418 & 5.3\% \\
T6 & 265 & 87.0\% & 12,405 & 89.6\% \\
T4 & 1.5 & 0.2\% & 29 & 0.1\% \\
T7 & 5.2 & 5.2\% & 106 & 2.3\% \\
T1 & 11.3 & 3.8\% & 237 & 1.7\% \\
T2, T5, T8, T9 & 1.4 to 1.5 each & 0.5\% each & 32 to 36 each & 0.2 to 0.3\% each \\
\bottomrule\end{longtable}}
{\itshape Table 5.2.2: expt\_b2 latencies and share of total time per transaction (64-client values are flagged diagnostic observations).}\par
Looking at the proportion of time spent on each transaction type, T6 now became the highest, at about 90\% of all time.\par
\needspace{4\baselineskip}\subsection{\texttt{expt\_\allowbreak{}b3}}
T6 aggregates, per merchant, the total number of orders and revenue amount in the last 30 days (via \texttt{created\_\allowbreak{}at} column), and returns the top 20. Under the vanilla baseline definition, there is no index on \texttt{created\_\allowbreak{}at}, hence Postgres needs to perform sequential table scan.\par
Postgres stores a table in 8KB pages, each holding the data tuples. When it needs a row, it reads the entire page, so the cost of a query is mostly the number of pages it reads. From \texttt{run.sh}, orders timestamp are generated in the following fashion:\par
\texttt{created\_\allowbreak{}at} \texttt{= current\_\allowbreak{}timestamp - (id * 31) mod 365 days}\par
so qualifying recent rows are spread across many heap pages. An ordinary B-Tree index can quickly list all the rows in the last 30 days, but it stores only the indexed column, and not the other information that T6 requires. To get those, Postgres still has to read the rows, which can require visiting a large fraction of the heap pages.\par
\needspace{4\baselineskip}\subsubsection*{Key Changes}
Create a covering index, \texttt{orders(created\_\allowbreak{}at) INCLUDE (merchant\_\allowbreak{}id, total\_\allowbreak{}amount)} that is sorted by date, and each entry carries the other two columns T6 needs. This enables an index-only scan when the visibility map permits heap access to be skipped \citep{pg-index-only}.\par
Perform \texttt{VACUUM} to mark table pages all-visible in the visibility map, enabling index-only scans to skip the table.\par
\needspace{4\baselineskip}\subsubsection*{Benchmark Results and Analysis}
\needspace{10\baselineskip}
{\fontsize{10.5}{12.5}\selectfont\setlength{\tabcolsep}{2mm}\renewcommand{\arraystretch}{1.15}
\begin{longtable}{>{\raggedright\arraybackslash}p{30.0mm}>{\raggedright\arraybackslash}p{30.0mm}>{\raggedright\arraybackslash}p{30.0mm}>{\raggedright\arraybackslash}p{30.0mm}>{\raggedright\arraybackslash}p{30.0mm}}\toprule
 & \textbf{TPS} & \textbf{p50 Latency (ms)} & \textbf{p95 Latency (ms)} & \textbf{p99 Latency (ms)} \\
\midrule\endhead
T6 & 4.61 \ensuremath{\rightarrow} 24.48 & 12,443 \ensuremath{\rightarrow} 1,213 & 13,782 \ensuremath{\rightarrow} 1,725 & 14,275 \ensuremath{\rightarrow} 2,038 \\
T3 & 2.37 \ensuremath{\rightarrow} 12.52 & 1,381 \ensuremath{\rightarrow} 607 & 2,139 \ensuremath{\rightarrow} 1,294 & 2,422 \ensuremath{\rightarrow} 2,109 \\
Overall & 46.04 \ensuremath{\rightarrow} 249.24 & 68 \ensuremath{\rightarrow} 35 & 12,444 \ensuremath{\rightarrow} 1,231 & 13,510 \ensuremath{\rightarrow} 1,619 \\
\bottomrule\end{longtable}}
{\itshape Table 5.3.1: Comparing expt\_b3 against expt\_b2 on T3, T6 at 64 clients  --  flagged raw observations; aborted phases exclude a clean comparison.}\par
b3's c1 phases complete at 97.932 TPS with no measured failures. c4-c64 measurements abort, and c8-c64 warmups also abort. Their raw latency/TPS improvements are retained as diagnostic observations; the c64 5-fold ratio does not establish clean performance gain.\par
\needspace{4\baselineskip}\subsection{\texttt{expt\_\allowbreak{}b4}}
The saved logs show worker serialization cancellations and Citus shard-commit errors that pgbench treats as client aborts (Figure 5.4.1). b4 tests READ ONLY together with coordinator-resolved T6 timestamps and plan capture; its node allocation also differs from b3. The resulting reduction in aborts is an observation, not an isolated READ ONLY causal test or a general concurrency fix.\par
\needspace{10\baselineskip}
{\fontsize{10.5}{12.5}\selectfont\setlength{\tabcolsep}{2mm}\renewcommand{\arraystretch}{1.15}
\begin{longtable}{>{\raggedright\arraybackslash}p{166.0mm}}\toprule
\texttt{WARNING:  failed to commit transaction on xcnf11:5446} \texttt{pgbench: error: client 2 script 0 aborted in command 153 query 0: ERROR:  could not commit transaction for shard 102081 on at least one active node} \\
\midrule\endhead
\bottomrule\end{longtable}}
{\itshape Figure 5.4.1: Log snippets of serialization error}\par
\needspace{4\baselineskip}\subsubsection*{Key Changes}
T3, T4 and T6 begin with \texttt{BEGIN READ ONLY} instead of \texttt{BEGIN}.\par
T6 reads the current time once on the coordinator (\texttt{LOCALTIMESTAMP}) and passes the window to its query as two literal timestamps.\par
\texttt{run.sh} saves \texttt{EXPLAIN (ANALYZE, VERBOSE, BUFFERS)} output for T6, T3 and T5.\par
\needspace{4\baselineskip}\subsubsection*{Benchmark Results and Analysis}
\needspace{10\baselineskip}
{\fontsize{10.5}{12.5}\selectfont\setlength{\tabcolsep}{2mm}\renewcommand{\arraystretch}{1.15}
\begin{longtable}{>{\raggedright\arraybackslash}p{30.0mm}>{\raggedright\arraybackslash}p{30.0mm}>{\raggedright\arraybackslash}p{30.0mm}>{\raggedright\arraybackslash}p{30.0mm}>{\raggedright\arraybackslash}p{30.0mm}}\toprule
 & \textbf{TPS} & \textbf{p50 Latency (ms)} & \textbf{p95 Latency (ms)} & \textbf{p99 Latency (ms)} \\
\midrule\endhead
T6 & 24.48 \ensuremath{\rightarrow} 31.17 & 1,213 \ensuremath{\rightarrow} 1,242 & 1,725 \ensuremath{\rightarrow} 1,758 & 2,038 \ensuremath{\rightarrow} 2,075 \\
Overall & 249.24 \ensuremath{\rightarrow} 316.00 & 35 \ensuremath{\rightarrow} 37 & 1,231 \ensuremath{\rightarrow} 1,242 & 1,619 \ensuremath{\rightarrow} 1,656 \\
\bottomrule\end{longtable}}
{\itshape Table 5.4.1: Comparing expt\_b4 against expt\_b3 on T6 at 64 clients  --  flagged raw observations; aborted phases exclude a clean comparison.}\par
Measured client-abort messages total 41 in b3 and two in b4; b4 still exits 2 at c32/c64. b4 has complete zero-failure phases at 1/4/8/16 clients. Differences in SQL and allocation prevent attribution to READ ONLY alone. Recorded worker cancellations and retry counts remain useful diagnostics; do not rank the aborted c64 runs as clean.\par
Separately, analysis of captured query \texttt{EXPLAIN} plans re-confirms \texttt{expt\_\allowbreak{}b3} changes that T6 is fully on index-only scans.\par
\needspace{10\baselineskip}
{\fontsize{10.5}{12.5}\selectfont\setlength{\tabcolsep}{2mm}\renewcommand{\arraystretch}{1.15}
\begin{longtable}{>{\raggedright\arraybackslash}p{166.0mm}}\toprule
\texttt{->  Index Only Scan using idx\_\allowbreak{}orders\_\allowbreak{}created\_\allowbreak{}at\_\allowbreak{}covering\_\allowbreak{}102104 on public.orders\_\allowbreak{}102104 orders  (cost=0.42..794.50 rows=13104 width=14)} \texttt{  Output: created\_\allowbreak{}at, merchant\_\allowbreak{}id, total\_\allowbreak{}amount} \texttt{  Index Cond: ((orders.created\_\allowbreak{}at >= '2026-09-11 12:42:13.736861'::timestamp without time zone) AND (orders.created\_\allowbreak{}at < '2026-10-11 12:42:13.736861'::timestamp without time zone))} \\
\midrule\endhead
\bottomrule\end{longtable}}
{\itshape Figure 5.4.2: Snippet of T6 query plan}\par
\needspace{4\baselineskip}\subsection{\texttt{expt\_\allowbreak{}b5}}
The same query plan generated in \texttt{expt\_\allowbreak{}b4} discovered that T3 is currently reading all the merchant's orders from the table via \texttt{merchant\_\allowbreak{}id} and then discarding rows not related to \texttt{customer\_\allowbreak{}id}.\par
\needspace{10\baselineskip}
{\fontsize{10.5}{12.5}\selectfont\setlength{\tabcolsep}{2mm}\renewcommand{\arraystretch}{1.15}
\begin{longtable}{>{\raggedright\arraybackslash}p{166.0mm}}\toprule
\texttt{->  Bitmap Heap Scan on public.orders\_\allowbreak{}102124 o  (cost=14.12..1779.77 rows=1 width=32) (actual time=8.631..8.632 rows=0.00 loops=1)} \texttt{      Output: o.order\_\allowbreak{}id, o.merchant\_\allowbreak{}id, o.customer\_\allowbreak{}id, o.order\_\allowbreak{}status, o.total\_\allowbreak{}amount, o.created\_\allowbreak{}at} \texttt{      Recheck Cond: (o.merchant\_\allowbreak{}id = 1)} \texttt{      Filter: (o.customer\_\allowbreak{}id = 1)} \texttt{      Rows Removed by Filter: 727} \texttt{      Heap Blocks: exact=571} \texttt{      Buffers: shared hit=578} \texttt{      ->  Bitmap Index Scan on idx\_\allowbreak{}orders\_\allowbreak{}merchant\_\allowbreak{}id\_\allowbreak{}102124  (cost=0.00..14.12 rows=760 width=0) (actual time=0.532..0.533 rows=727.00 loops=1)} \texttt{            Index Cond: (o.merchant\_\allowbreak{}id = 1)} \texttt{            Index Searches: 1} \texttt{            Buffers: shared hit=7} \\
\midrule\endhead
\bottomrule\end{longtable}}
{\itshape Figure 5.5.1: Snippet of T3 query plan in expt\_b4}\par
\needspace{4\baselineskip}\subsubsection*{Key Changes}
\texttt{CREATE INDEX ... ON orders(merchant\_\allowbreak{}id, customer\_\allowbreak{}id)}\par
\needspace{4\baselineskip}\subsubsection*{Benchmark Results and Analysis}
\needspace{10\baselineskip}
{\fontsize{10.5}{12.5}\selectfont\setlength{\tabcolsep}{2mm}\renewcommand{\arraystretch}{1.15}
\begin{longtable}{>{\raggedright\arraybackslash}p{30.0mm}>{\raggedright\arraybackslash}p{30.0mm}>{\raggedright\arraybackslash}p{30.0mm}>{\raggedright\arraybackslash}p{30.0mm}>{\raggedright\arraybackslash}p{30.0mm}}\toprule
 & \textbf{TPS} & \textbf{p50 Latency (ms)} & \textbf{p95 Latency (ms)} & \textbf{p99 Latency (ms)} \\
\midrule\endhead
T3 & 15.83 \ensuremath{\rightarrow} 15.80 & 707 \ensuremath{\rightarrow} 620 & 990 \ensuremath{\rightarrow} 924 & 1,244 \ensuremath{\rightarrow} 1,111 \\
Overall & 316.00 \ensuremath{\rightarrow} 315.12 & 37 \ensuremath{\rightarrow} 41 & 1,242 \ensuremath{\rightarrow} 1,263 & 1,656 \ensuremath{\rightarrow} 1,650 \\
\bottomrule\end{longtable}}
{\itshape Table 5.5.1: Comparing expt\_b5 against expt\_b4 on T3 at 64 clients  --  flagged raw observations; aborted phases exclude a clean comparison.}\par
\needspace{10\baselineskip}
{\fontsize{10.5}{12.5}\selectfont\setlength{\tabcolsep}{2mm}\renewcommand{\arraystretch}{1.15}
\begin{longtable}{>{\raggedright\arraybackslash}p{166.0mm}}\toprule
\texttt{->  Index Scan using idx\_\allowbreak{}orders\_\allowbreak{}merchant\_\allowbreak{}customer\_\allowbreak{}102117 on public.orders\_\allowbreak{}102117 o  ...} \texttt{Planning Time: 0.633 ms} \texttt{Execution Time: 0.306 ms} \\
\midrule\endhead
\bottomrule\end{longtable}}
{\itshape Figure 5.5.2: Snippet of T3 query plan in expt\_b5}\par
The captured T3 shard plan is faster after the merchant/customer index, while T6 remains prominent in recorded time shares. b4 and b5 c64 measurements both abort (b5's warmup also aborts), so their similar raw TPS is not an eligible matched throughput comparison.\par
\needspace{4\baselineskip}\subsection{\texttt{expt\_\allowbreak{}b6}}
One approach to reduce the amount of data read by analytics queries such as T6 is pre-aggregating and summarising raw tables into rollup tables. In our case, the order counts and revenue per day before the load date can be computed once and reused by every T6.\par
\needspace{4\baselineskip}\subsubsection*{Key Changes}
A new Citus Reference table, \texttt{merchant\_\allowbreak{}daily\_\allowbreak{}sales(sales\_\allowbreak{}day, merchant\_\allowbreak{}id, order\_\allowbreak{}count, revenue)}, with primary key \texttt{(sales\_\allowbreak{}day, merchant\_\allowbreak{}id)} and the two totals included in the index so it can be read by index-only scans. Reference tables are normally expensive to write, because every write goes to every node, but for this benchmarking experiment, the tradeoff could be worth it since this table is fixed throughout the experiment run.\par
\texttt{merchant\_\allowbreak{}daily\_\allowbreak{}sales\_\allowbreak{}cutoff(cutoff\_\allowbreak{}day)}: a one-row plain table on the coordinator that records the first day not in the rollup.\par
At the end of every load, \texttt{run.sh} records the cutoff, fills the rollup table with one \texttt{INSERT ... SELECT ... GROUP BY} over \texttt{orders}, and runs \texttt{VACUUM} on it.\par
T6 splits its query window into three parts: the partial first day and the part since the cutoff, both read from orders, plus whole days read from the rollup. This static rollup relies on the supplied workload: T1 inserts current-time orders, while no transaction backdates orders or changes historical order amounts. It would require maintenance if those assumptions changed.\par
The rollup stores daily totals rather than results for a fixed \texttt{REPORT\_DAYS} value. Each T6 call selects the relevant days from the same rollup and reads uncovered periods from the original orders table, so changing \texttt{REPORT\_DAYS} does not require rebuilding the rollup.\par
\needspace{4\baselineskip}\subsubsection*{Benchmark Results and Analysis}
\needspace{10\baselineskip}
{\fontsize{10.5}{12.5}\selectfont\setlength{\tabcolsep}{2mm}\renewcommand{\arraystretch}{1.15}
\begin{longtable}{>{\raggedright\arraybackslash}p{30.0mm}>{\raggedright\arraybackslash}p{30.0mm}>{\raggedright\arraybackslash}p{30.0mm}>{\raggedright\arraybackslash}p{30.0mm}>{\raggedright\arraybackslash}p{30.0mm}}\toprule
 & \textbf{TPS} & \textbf{p50 Latency (ms)} & \textbf{p95 Latency (ms)} & \textbf{p99 Latency (ms)} \\
\midrule\endhead
T6 & 31.05 \ensuremath{\rightarrow} 47.90 & 1,266 \ensuremath{\rightarrow} 870 & 1,757 \ensuremath{\rightarrow} 1,077 & 2,022 \ensuremath{\rightarrow} 1,174 \\
T3 & 15.80 \ensuremath{\rightarrow} 24.17 & 620 \ensuremath{\rightarrow} 161 & 924 \ensuremath{\rightarrow} 226 & 1,111 \ensuremath{\rightarrow} 261 \\
Overall & 315.12 \ensuremath{\rightarrow} 483.53 & 41 \ensuremath{\rightarrow} 39 & 1,263 \ensuremath{\rightarrow} 869 & 1,650 \ensuremath{\rightarrow} 1,031 \\
\bottomrule\end{longtable}}
{\itshape Table 5.6.1: Comparing expt\_b6 against expt\_b5 on T6 at 64 clients  --  flagged raw observations; aborted phases exclude a clean comparison.}\par
Because this rewrite could silently change T6's answer, we added additional verification. After the load and after every measured run, we run the original T6 query and the new one with the same window and compare all 1,000 per-merchant totals in both directions with \texttt{EXCEPT}. We compared the full totals rather than the top 20, because merchants tied for 20th place could legitimately appear in either order. The validation check found no differences between the original and optimized T6 queries. The final separate correctness job-1031186 found no negative inventory, order-total mismatches or duplicate inventory/order-item IDs before and after T7 and after a 300-second run with 4 clients. Two isolated T7 transfers also passed, with total stock remaining at 1,773. This separate job completed both phases with zero failed records and four retry attempts; it does not retroactively validate earlier aborted benchmarks.\par
\needspace{4\baselineskip}\subsubsection*{Matched baseline comparison and separate correctness run}
We compared \texttt{expt\_b0}/job 1029473 and \texttt{expt\_b6}/job 1030000  on the same nodes xcnf0--4, using scale 1, seed 42, 60-second warmups and 600-second measurements. The original b0 run used a different node allocation, so we use this new matched baseline for comparison.\par
\needspace{10\baselineskip}
{\fontsize{10.5}{12.5}\selectfont\setlength{\tabcolsep}{2mm}\renewcommand{\arraystretch}{1.15}
\begin{longtable}{>{\raggedright\arraybackslash}p{30.0mm}>{\raggedright\arraybackslash}p{30.0mm}>{\raggedright\arraybackslash}p{30.0mm}>{\raggedright\arraybackslash}p{30.0mm}>{\raggedright\arraybackslash}p{30.0mm}}\toprule
 & \textbf{TPS} & \textbf{p50 Latency (ms)} & \textbf{p95 Latency (ms)} & \textbf{p99 Latency (ms)} \\
\midrule\endhead
b0 & 5.190 & 16.2 & 464.0 & 2679.0 \\
b6 & 108.250 & 4.4 & 56.1 & 61.9 \\
\bottomrule\end{longtable}}
{\itshape Table 5.6.2: Matched expt\_b0 and expt\_b6 at 1 client (seed 42). Both runs completed without failed records.}\par
At 1 client, b6 achieved 20.86 times the baseline throughput, with much lower tail latency. At 4 and 16 clients, b6 completed at 394.380 and 554.982 TPS, while b0 either aborted or recorded final serialization failures. b6 also aborted at 8 clients, so these cases cannot be used for direct speedup comparisons.\par
We selected b6 as the final B design because it combines the earlier placement and indexing optimizations with T6 pre-aggregation, achieving a 20.86\ensuremath{\times} throughput improvement in the matched one-client comparison while passing the separate correctness checks.\par

\paragraph{Seven-day validation.}
A separate scale-1 diagnostic (\texttt{expt\_b6}, job 1041191) used \texttt{REPORT\_DAYS=7}, one client, a 30-second measurement and no warmup. Before and after the workload, the original and optimized T6 queries returned identical order counts and revenue for all 1,000 merchants, with zero differences in both directions. All 3,988 scripts, including 397 T6 executions, completed without recorded failures or retries. This supplements the 30-day evidence; it does not test per-call window changes or zero-day/boundary cases. Its timing is not a matched performance comparison with the full 30-day runs.\par
